\section{Trace of an Endomorphism of Finite Rank}

In this subsection, A denotes a commutative ring. Let E be a projective A-module. Then \(\operatorname{End}_{\mathrm{A}}^{\mathrm{f}}(\mathrm{E})\) is an A-submodule of \(\operatorname{End}_{\mathrm{A}}(\mathrm{E})\) , and the canonical mapping \(E^{*}\otimes_{A}E\to\operatorname{End}_{\mathrm{A}}(\mathrm{E})\) defines an isomorphism \(\theta_{E}\) of A-modules from \(E^{*}\otimes_{A}E\) to \(\operatorname{End}_{\mathrm{A}}^{\mathrm{f}}(\mathrm{E})\) (Lemma 1). Consider the canonical linear form \(\tau:E^{*}\otimes_{A}E\to A\) (II, §4, No.~3, p.~273) characterized by the formula \(\tau(x^{*}\otimes x)=\langle x,x^{*}\rangle\) . By composing it with the isomorphism \(\theta_{E}^{-1}\) , we deduce a linear form \(\operatorname{Tr}: \operatorname{End}_{\mathrm{A}}^{\mathrm{f}}(\mathrm{E})\to\mathrm{A}\) , called the trace form. When the A-module E is finitely generated, we recover the definition of II, §4, No.~3, p.~273.

PROPOSITION 1. --- Let E be a free A-module. Let \((e_{i})_{i\in\mathrm{I}}\) be a basis of E, and let \((e_{i}^{*})_{i\in\mathrm{I}}\) be its dual basis. Let \(u\in\operatorname{End}_{\mathrm{A}}^{\mathrm{f}}(\mathrm{E})\) . The family \((\langle u(e_{i}),e_{i}^{*}\rangle)_{i\in\mathrm{I}}\) has finite support, and its sum is equal to \(\operatorname{Tr}(u)\) .

It suffices to treat the case when \(u\) is of the form \(\theta_{\mathrm{E}}(x^{*} \otimes x)\) with \(x \in \mathrm{E}\) and \(x^{*} \in \mathrm{E}^{*}\) . The family \((\langle x, e_{i}^{*} \rangle)_{i \in \mathrm{I}}\) then has finite support, and we have \(x = \sum_{i \in \mathrm{I}} \langle x, e_{i}^{*} \rangle e_{i}\) . Consequently, the family \((\langle x, e_{i}^{*} \rangle \langle e_{i}, x^{*} \rangle)_{i \in \mathrm{I}}\) also has finite support, and we have \(\langle x, x^{*} \rangle = \sum_{i \in \mathrm{I}} \langle x, e_{i}^{*} \rangle \langle e_{i}, x^{*} \rangle\) . Now, we have \(\langle u(e_{i}), e_{i}^{*} \rangle = \langle x, e_{i}^{*} \rangle \langle e_{i}, x^{*} \rangle\) for every \(i \in \mathrm{I}\) . This proves the proposition.

PROPOSITION 2. --- Let E, F be projective A-modules. Let \(u \in \operatorname{Hom}_{\mathrm{A}}^{\mathrm{f}}(\mathrm{E}, \mathrm{F})\) and \(v \in \operatorname{Hom}_{\mathrm{A}}(\mathrm{F}, \mathrm{E})\) . We have the relation

\[
\operatorname{Tr} (v \circ u) = \operatorname{Tr} (u \circ v).\tag{3}
\]

It suffices to prove the proposition when \(u\) is of the form \(\theta(x^{*} \otimes y)\) with \(x^{*} \in \mathrm{E}^{*}\) and \(y \in \mathrm{F}\) . In this case, we have

\[
v \circ u = \theta_ {\mathrm{E}} (x ^ {*} \otimes v (y)) \quad \text { and } \quad u \circ v = \theta_ {\mathrm{F}} (^ {t} v (x ^ {*}) \otimes y)  ,
\]

and therefore

\[
\operatorname{Tr} (v \circ u) = \langle v (y), x ^ {*} \rangle = \langle y, ^ {t} v (x ^ {*}) \rangle = \operatorname{Tr} (u \circ v).
\]

COROLLARY. --- Let E be a projective A-module, u be an element of \(\operatorname{End}_{\mathrm{A}}^{f}(\mathrm{E})\) , and F be a projective A-submodule of E containing Im u. Denote by \(u_{F}\) the endomorphism of F induced by u. We have

\[
\operatorname{Tr} (u) = \operatorname{Tr} (u _ {\mathrm{F}}).\tag{4}
\]

Denote by i the canonical injection of F into E and by \(v: E \rightarrow F\) the homomorphism deduced from u. We have \(u_{F} = v \circ i\) and \(u = i \circ v\) . The corollary follows.

Let E be a projective A-module and \(u \in \operatorname{End}_{\mathrm{A}}^{\mathrm{f}}(\mathrm{E})\) . For every natural number p, the A-module \(\wedge^{p}E\) is projective (III, §7, No.~8, p.~519, Corollary 2), and the endomorphism \(\wedge^{p}u\) belongs to \(\operatorname{End}_{\mathrm{A}}^{\mathrm{f}}(\wedge^{p}\mathrm{E})\) (III, §7, No.~3, p.~511, Proposition 6) and is zero for p sufficiently large. The set \(1_{\mathrm{E}} + \operatorname{End}_{\mathrm{A}}^{\mathrm{f}}(\mathrm{E})\) is stable under composition. We define a mapping det from \(1_{\mathrm{E}} + \operatorname{End}_{\mathrm{A}}^{\mathrm{f}}(\mathrm{E})\) to A by setting

\[
\det (1 _ {\mathrm{E}} + u) = \sum_ {p \geqslant 0} \operatorname{Tr} \wedge^ {p} u
\]

for \(u\in \mathrm{End}_{\mathrm{A}}^{\mathrm{f}}(\mathrm{E})\)

If E is free and finite-dimensional, then this definition agrees with that of III, §8, No.~1, p.~522, by the corollary of III, §8, No.~5, p.~530.

PROPOSITION 3. --- Let E be a projective A-module.

\begin{enumerate}
\def\labelenumi{\alph{enumi})}
\tightlist
\item
  Let \(u \in \mathrm{End}_{\mathrm{A}}^{\mathrm{f}}(\mathrm{E})\) . Let \(\mathrm{F}\) be a projective A-submodule of \(\mathrm{E}\) containing \(\operatorname{Im} u\) , and let \(u_{\mathrm{F}}\) be the endomorphism of \(\mathrm{F}\) induced by \(u\) . We have
\end{enumerate}

\[
\det (1 _ {\mathrm{E}} + u) = \det (1 _ {\mathrm{F}} + u _ {\mathrm{F}}).\tag{5}
\]

\begin{enumerate}
\def\labelenumi{\alph{enumi})}
\setcounter{enumi}{1}
\tightlist
\item
  Let \(u\) and \(v\) be two elements of \(\mathrm{End}_{\mathrm{A}}^{\mathrm{f}}(\mathrm{E})\) . We have
\end{enumerate}

\[
\det \bigl ((1 _ {\mathrm{E}} + u) \circ (1 _ {\mathrm{E}} + v) \bigr) = \det (1 _ {\mathrm{E}} + u) \det (1 _ {\mathrm{E}} + v).\tag{6}
\]

Let us prove a). For every integer \(p \geqslant 0\) , the projective A-module \(\wedge^p\mathrm{F}\) can be identified with a submodule of \(\wedge^p\mathrm{E}\) (III, §7, No.~9, p.~520, Corollary). The image of \(\wedge^p u\) is contained in \(\wedge^p\mathrm{F}\) , and the endomorphism of \(\wedge^p\mathrm{F}\) induced by \(\wedge^{p}u\) is equal to \(\wedge^{p}u_{F}\) . We consequently have \(\operatorname{Tr}(\wedge^{p}u)=\operatorname{Tr}(\wedge^{p}u_{F})\) by the corollary of Proposition 2, and therefore a).

Let us prove b). Let G be an A-module such that the A-module \(L = E \oplus G\) is free. Denote by \(u'\) and \(v'\) the endomorphisms \(u \oplus 0_{G}\) and \(v \oplus 0_{G}\) of L. By a), we have the relations \(\det(1_{L} + u') = \det(1_{E} + u)\) , \(\det(1_{L} + v') = \det(1_{E} + v)\) , and

\[
\det (1 _ {\mathrm{L}} + u ^ {\prime} + v ^ {\prime} + u ^ {\prime} \circ v ^ {\prime}) = \det (1 _ {\mathrm{E}} + u + v + u \circ v).
\]

It therefore suffices to prove assertion b) when the A-module E is free. There then exists a finitely generated free submodule F of E that contains the image of u and that of v. Set \(w = u + v + u \circ v\) . The image of w is contained in F, and we have \(w_{F} = u_{F} + v_{F} + u_{F} \circ v_{F}\) . Therefore, by (5), we have \(\det(1_{E} + u) = \det(1_{F} + u_{F})\) , \(\det(1_{E} + v) = \det(1_{F} + v_{F})\) , and

\[
\begin{array}{c} \det \bigl ((1 _ {\mathrm{E}} + u) \circ (1 _ {\mathrm{E}} + v) \bigr) \\ = \det (1 _ {\mathrm{E}} + w) = \det (1 _ {\mathrm{F}} + w _ {\mathrm{F}}) = \det \bigl ((1 _ {\mathrm{F}} + u _ {\mathrm{F}}) \circ (1 _ {\mathrm{F}} + v _ {\mathrm{F}}) \bigr). \end{array}
\]

Since F is a finitely generated free A-module, we have

\[
\det \left(\left(1 _ {\mathrm{F}} + u _ {\mathrm{F}}\right) \circ \left(1 _ {\mathrm{F}} + v _ {\mathrm{F}}\right)\right) = \det \left(1 _ {\mathrm{F}} + u _ {\mathrm{F}}\right) \det \left(1 _ {\mathrm{F}} + v _ {\mathrm{F}}\right) = \det \left(1 _ {\mathrm{E}} + u\right) \det \left(1 _ {\mathrm{E}} + v\right).
\]

